\documentclass[11pt]{article}
\usepackage{docstyle}
% Lesson 4 lost-mark points (from the material plus teaching observation), numbered once
\setcounter{lostmark}{0}
\lostmark{plusc}{Forgetting $+c$, or finding $c$ with the wrong point (use the point the question gives).}
\lostmark{newpower}{Dividing by the old power: $x^{3}\to\dfrac{x^{4}}{4}$, divide by the \kw{new} power.}
\lostmark{value}{Rate of change: giving the value of the quantity instead of its \tderiv.}
\lostmark{meaning}{No meaning or units: say ``decreasing at 528 people per hour'', not just $-528$.}
\lostmark{wrongzero}{Kinematics: maximum height or ``at rest'' needs $v=0$, not $a=0$ or $s=0$.}
\lostmark{distance}{Distance is not displacement: if the object turns round, add the distance of each part.}

\newlength\figunit
\begin{document}
\handouthead{Lesson 4 \textperiodcentered\ Anti-differentiation, Rates of Change and Kinematics}{NCEA Level 2 Mathematics \textperiodcentered\ AS 91262 Apply calculus methods in solving problems \textperiodcentered\ Handout}

\section{The one idea}\label{idea}
\begin{keybox}
\kw{Anti-differentiation undoes differentiation.} Raise the power by 1, divide by the new power, add $+c$.
Then use a known point to find $c$.
\end{keybox}
A \tderiv is a \kw{rate of change}. In motion, velocity is the rate of change of displacement,
and acceleration is the rate of change of velocity.

\section{Anti-differentiating}\label{anti}
\begin{keybox}
\[x^{n}\ \longrightarrow\ \dfrac{x^{n+1}}{n+1}+c\qquad\qquad k\ \longrightarrow\ kx+c\]
\end{keybox}
\begin{tabular}{@{}llll@{}}
\toprule
$f'(x)$ term & raise the power & divide by the new power & result\\
\midrule
$6x^{2}$ & $6x^{3}$ & $\dfrac{6x^{3}}{3}$ & $2x^{3}$\\[2mm]
$-4x$ & $-4x^{2}$ & $\dfrac{-4x^{2}}{2}$ & $-2x^{2}$\\[2mm]
$5$ & $5x$ & $\dfrac{5x}{1}$ & $5x$\\
\bottomrule
\end{tabular}

\medskip
\kw{Why $+c$?} Differentiating removes a constant, so the \tantider can have any constant.
\kw{Check}: differentiate your answer; you must get back $f'(x)$.

\begin{example}{Worked example 1}
Find $f(x)$ if $f'(x)=6x^{2}-4x+5$.
\begin{steps}
\item Term by term (table above).
\item $f(x)=\mathbf{2x^{3}-2x^{2}+5x+c}$.\quad Check: $6x^{2}-4x+5$.
\end{steps}
\end{example}

\section{Finding $c$}\label{findc}
\begin{example}{Worked example 2}
$f'(x)=4x-3$ and the curve passes through $(2,\,7)$. Find $f(x)$.
\begin{steps}
\item $f(x)=2x^{2}-3x+c$
\item Substitute the point: $7=2(2)^{2}-3(2)+c=2+c$, so $c=5$.
\item $f(x)=\mathbf{2x^{2}-3x+5}$
\end{steps}
\end{example}

\section{Rates of change}\label{rates}
\subsection{Write it in this order}
\begin{enumerate}[leftmargin=7mm, itemsep=1pt]
\item Differentiate: the \tderiv is the rate.
\item Substitute the given time (if the question gives a \emph{value of the quantity}, solve for the time first).
\item Answer with units ``per \dots'' and the meaning: positive $=$ increasing, negative $=$ decreasing.
\end{enumerate}

\begin{example}{Worked example 3}
The water in a tank is $V=2000+30t-0.5t^{2}$ litres after $t$ minutes. Find the rate of change at $t=10$ and at $t=40$.
\begin{steps}
\item $\dd{V}{t}=30-t$
\item $t=10$: $20$, so the volume is \kw{increasing at $20$ litres per minute}.
\item $t=40$: $-10$, so the volume is \kw{decreasing at $10$ litres per minute}.
\end{steps}
\end{example}

\begin{example}{Worked example 4 \textperiodcentered\ given the rate, find the time}
A population is $P=500+40t-t^{2}$ after $t$ years. When is it growing at $10$ per year?
\begin{steps}
\item $\dd{P}{t}=40-2t=10$
\item $2t=30$, so $\mathbf{t=15}$ \textbf{years}.
\end{steps}
\end{example}

\textbf{Sentence frame:} ``At $t=\dots$ the \dots\ is increasing/decreasing at a rate of \dots\ per \dots''

\section{Kinematics: displacement, velocity, acceleration}\label{kin}
\begin{center}
\setlength\figunit{9mm}
% displacement -> velocity -> acceleration: differentiate to the right, anti-differentiate to the left
\begin{tikzpicture}[x=\figunit, y=\figunit]
  \foreach \x/\t/\n in {0/{\tdisp}/s, 4.6/{\tvel}/v, 9.2/{\tacc}/a}{
    \draw[line width=0.8pt, fill=boxfill, rounded corners=2pt] (\x-1.65,-0.65) rectangle (\x+1.65,0.65);
    \node at (\x,0.15) {$\n$};
    \node at (\x,-0.3) {\t};}
  \foreach \x in {0,4.6}{
    \draw[-{Stealth[length=2mm]}, line width=0.8pt] (\x+1.75,0.35) -- (\x+2.85,0.35);
    \draw[-{Stealth[length=2mm]}, line width=0.8pt] (\x+2.85,-0.35) -- (\x+1.75,-0.35);}
  \node[above] at (4.6,0.75) {differentiate $\longrightarrow$};
  \node[below] at (4.6,-0.75) {$\longleftarrow$ anti-differentiate, then find $c$};
\end{tikzpicture}

\end{center}

\begin{tabular}{@{}p{0.42\linewidth}p{0.52\linewidth}@{}}
\toprule
words in the question & maths\\
\midrule
at rest, stops, changes direction & $v=0$\\
maximum height & $v=0$ (then find $s$)\\
starts from the origin / the start line & $s=0$ when $t=0$\\
starts from rest & $v=0$ when $t=0$\\
initial velocity $u$ & $v=u$ when $t=0$\\
constant acceleration $-9.8$ (falling) & $a=-9.8\mpss$\\
\bottomrule
\end{tabular}

\begin{example}{Worked example 5 \textperiodcentered\ differentiate}
A particle moves with $s=t^{3}-6t^{2}+9t$ metres. Find (a) $v$ at $t=2$, (b) when it is at rest, (c) $a$ at $t=3$.
\begin{steps}
\item $v=\dd{s}{t}=3t^{2}-12t+9$;\quad (a) $v(2)=12-24+9=\mathbf{-3\mps}$
\item (b) $3(t-1)(t-3)=0$, so \textbf{at rest at $t=1$ s and $t=3$ s}.
\item $a=\dd{v}{t}=6t-12$;\quad (c) $a(3)=\mathbf{6\mpss}$
\end{steps}
\end{example}

\begin{example}{Worked example 6 \textperiodcentered\ anti-differentiate (Merit)}
A ball is thrown up from $1.2$ m above the ground at $14\mps$. Its acceleration is $-9.8\mpss$. Find its maximum height.
\begin{steps}
\item $v=-9.8t+c$; $v(0)=14$, so $v=-9.8t+14$.
\item $s=-4.9t^{2}+14t+c$; $s(0)=1.2$, so $s=-4.9t^{2}+14t+1.2$.
\item Maximum height when $v=0$: $t=\dfrac{14}{9.8}=1.4286$ s.
\item $s=-4.9(1.4286)^{2}+14(1.4286)+1.2=\mathbf{11.2}$ \textbf{m}.
\end{steps}
\end{example}

\section{What each grade looks like here}\label{grades}
\begin{tabular}{@{}p{0.14\linewidth}p{0.80\linewidth}@{}}
\toprule
Achieved & anti-differentiate and find $c$; find a rate at a given time; find $v$ or $a$ by differentiating.\\
Merit & find the time from a given rate or value; go from $a$ to $v$ to $s$ with two constants; interpret the sign of a rate.\\
Excellence & a multi-step motion problem (e.g.\ find the cliff height from the landing speed, then the maximum height), fully explained.\\
\bottomrule
\end{tabular}

\section{Checking with the fx-9750GIII}\label{calc}
In \textsf{RUN-MATRIX}, \textsf{OPTN} $\to$ \textsf{CALC} also has $\int\mathrm{d}x$:
the definite integral of $v$ from $0$ to $3$ is the displacement over the first $3$ seconds.
Use it to check $s(3)-s(0)$; the working must still show $s(t)$ with its constant.

\section{Ways to lose marks}\label{lose}
From the material plus teaching observation.
\begin{enumerate}[leftmargin=7mm, itemsep=2pt]
\item[\textbf{\lmnum{plusc}}] \lmtxt{plusc}
\item[\textbf{\lmnum{newpower}}] \lmtxt{newpower}
\item[\textbf{\lmnum{value}}] \lmtxt{value}
\item[\textbf{\lmnum{meaning}}] \lmtxt{meaning}
\item[\textbf{\lmnum{wrongzero}}] \lmtxt{wrongzero}
\item[\textbf{\lmnum{distance}}] \lmtxt{distance}
\end{enumerate}
\end{document}
